\documentclass[11pt,letterpaper]{article}

\usepackage[margin=1in]{geometry}
\usepackage{amsmath,amssymb}
\usepackage{booktabs}
\usepackage{array}
\usepackage{graphicx}
\usepackage{enumitem}
\usepackage{tabularx}
\usepackage{titlesec}
\usepackage{xcolor}
\usepackage[hidelinks]{hyperref}

\definecolor{answerhl}{RGB}{198,239,206}
\definecolor{fillhl}{RGB}{255,235,156}
\newcommand{\hl}[1]{\colorbox{answerhl}{#1}}
\newcommand{\fb}[1]{\colorbox{fillhl}{\textit{#1}}}

\titleformat{\section}{\Large\bfseries}{\thesection}{1em}{}
\titleformat{\subsection}{\large\bfseries}{\thesubsection}{1em}{}
\setlength{\parskip}{0.5em}
\setlength{\parindent}{0pt}

\begin{document}

\begin{titlepage}
\centering
\vspace*{2cm}
{\Huge\bfseries\scshape Portland State \\[2pt] University \par}
\vspace{1.5cm}
{\large\scshape Department of \par}
{\large\scshape Electrical and Computer Engineering \par}
\vspace{1.5cm}
{\LARGE\bfseries ECE332 \par}
\vspace{0.4cm}
{\LARGE\scshape Lab 4: Patch Antenna \par}
\vfill
{\large \textbf{Prepared by:} \par}
{\large Bao Nguyen \par}
{\large Nick Stites \par}
\vspace{1cm}
{\large \textbf{Term:} Spring 2026 \par}
\vspace{0.4cm}
{\large \textbf{Date:} May 14, 2026 \par}
\vspace{2cm}
\end{titlepage}

\thispagestyle{empty}
\tableofcontents
\newpage

\section{Introduction}

This lab builds and characterizes a half-wave rectangular patch antenna on single-sided copper-clad FR4, tuned for the 2.4~GHz WiFi band ($f_c \approx 2.443$~GHz, the geometric mean of the 2.400--2.487~GHz channel range). The patch is fed at the edge by a 50~$\Omega$ microstrip transmission line made from copper tape, and a side-launch SMA connector ties the transmission line to the VNA. The patch length $L$ sets the resonant frequency through the half-wavelength-in-dielectric condition $f_c = c/(2L\sqrt{\varepsilon_r})$. The width $W$ controls the input impedance and the radiation pattern, and the inset distance $R$ tunes the feed-point impedance down toward 50~$\Omega$.

After construction we measure $S_{11}$ and $S_{22}$ on a VNA to verify the resonance lands inside the WiFi band, trim the patch length if needed, then cut an inset feed of length $R$ to drop the input impedance to 50~$\Omega$. The second half of the lab uses both antennas at once to measure $S_{21}$ and $S_{12}$ versus orientation angle, producing an azimuth pattern of the radiation lobe. A cross-polarization test ($90^\circ$ rotation of one antenna) verifies that patch antennas are linearly polarized in the far field.

\section{Antenna Design Calculations}

\subsection*{Patch length \texorpdfstring{$L$}{L} versus \texorpdfstring{$\varepsilon_r$}{epsilon\_r}}

The center frequency of the patch is set by the half-wavelength-in-dielectric condition,
\[
f_c \approx \frac{c}{2L\sqrt{\varepsilon_r}} \quad\Longrightarrow\quad L \approx \frac{c}{2 f_c \sqrt{\varepsilon_r}}.
\]
With $c = 2.998\times10^{8}$~m/s and $f_c = \sqrt{f_1 f_2} = \sqrt{2.400 \cdot 2.487} \approx 2.443$~GHz (the geometric mean of the WiFi band), sweeping $\varepsilon_r$ across the FR4 spec range gives:

\begin{center}
\begin{tabular}{cc}
\toprule
$\varepsilon_r$ & $L$ (mm) \\
\midrule
4.0 & 30.806 \\
4.1 & 30.428 \\
4.2 & 30.063 \\
4.3 & 29.711 \\
4.4 & 29.372 \\
\bottomrule
\end{tabular}
\end{center}

About the spread: a 10\% sweep in $\varepsilon_r$ (4.0 to 4.4) only moves $L$ by about 1.43~mm, which is roughly 4.7\% of $L$. The length scales as $1/\sqrt{\varepsilon_r}$, so the sensitivity is half the relative error in $\varepsilon_r$. For example a $\pm 0.1$ swing around $\varepsilon_r = 4.2$ shifts $L$ by about $\pm 0.36$~mm, which is right at the edge of what an Xacto knife can resolve. So basically picking the wrong $\varepsilon_r$ inside the FR4 spec range won't kill the resonance, it just nudges $f_c$ by a few tens of MHz and we trim by hand on the VNA anyway.

\textbf{Design choice:} use the middle value $L = 30$~mm for the initial cut, which corresponds to $\varepsilon_r \approx 4.2$.

\subsection*{Patch width \texorpdfstring{$W$}{W}}

The constraint from the lab is $W \geq L$, with larger $W$ dropping the input impedance and broadening the pattern. A square patch ($W = L \approx 30$~mm) sits around 300~$\Omega$. We picked $W = 67$~mm, which is the short edge of the FR4 board minus 2~mm of clearance margin. About 2.2$\times$ wider than $L$, which drops the edge-fed impedance well below the square-patch 300~$\Omega$ value and uses the full board width for the radiating aperture. $W$ runs along the short edge of the board so the transmission line and SMA connector fit along the long edge.

\subsection*{50~\texorpdfstring{$\Omega$}{ohm} transmission line width}

The 50~$\Omega$ microstrip trace width was computed using the Digi-Key PCB Trace Impedance Calculator\footnote{\url{https://www.digikey.com/en/resources/conversion-calculators/conversion-calculator-pcb-trace-impedance}} with the measured substrate parameters:

\begin{center}
\begin{tabular}{lc}
\toprule
$\varepsilon_r$         & 4.2 \\
Copper thickness $t$    & 0.04~mm \\
Substrate height $h$    & 1.40~mm \\
Target $Z_0$            & 50~$\Omega$ \\
\midrule
Trace width $w$         & 2.632~mm \\
\bottomrule
\end{tabular}
\end{center}

We used $w = 2.632$~mm for the transmission line cut. The slightly thinner-than-spec substrate ($h = 1.40$~mm measured vs $1.6$~mm nominal) and the copper-tape thickness ($t = 0.04$~mm) were both factored into the calculator instead of using the textbook approximation.

\subsection*{Wavelength and far-field boundary}

The free-space wavelength at $f_c = 2.443$~GHz is
\[
\lambda_0 = \frac{c}{f_c} = \frac{2.998\times 10^{8}}{2.443\times 10^{9}} \approx 12.28~\text{cm}.
\]
For an aperture antenna with largest dimension $D$, the far-field (Fraunhofer) region begins at
\[
r_{ff} = \frac{2 D^2}{\lambda_0}.
\]
With $D = W = 67$~mm, $r_{ff} \approx 2(0.067)^2/0.1228 \approx 7.31$~cm. So any cable-length separation (typically $\geq 30$~cm with the lab's blue SMAs) is comfortably in the far-field, and the azimuth pattern we measure will be the true radiation pattern rather than near-field coupling.

\section{VNA Setup}

The VNA was configured per the lab procedure with a 4-way quadrant split: top-left $S_{11}$ Log-Mag, top-right $S_{11}$ Smith (R+jX), bottom-left $S_{22}$ Log-Mag, bottom-right $S_{22}$ Smith (R+jX). Start frequency 1.5~GHz, stop frequency 2.7~GHz. The eCal module was used for calibration on both ports through the blue SMA cables. After the resonance was found in the WiFi band, the VNA was recalibrated over the narrower 2.35--2.55~GHz span.

\section{Build the Antennas}

Both antennas were built on single-sided copper-clad FR4. The patch and 50~$\Omega$ transmission line were cut from copper tape using an Xacto knife, with the transmission line meeting the patch at the center of the $W$-edge and running to the board edge where the side-launch SMA connector was soldered. The SMA center pin attaches to the transmission line and the body pins to the ground plane on the bare side of the board.

\subsection*{Photographs}

\fb{Insert photo of Bao's antenna with all five dimensions labeled in Sharpie on the board itself: $L$, $W$, $R$, transmission line length, transmission line width.}

\vspace{1em}
\fb{Insert photo of Nick's antenna with all five dimensions labeled in Sharpie on the board itself: $L$, $W$, $R$, transmission line length, transmission line width.}

\subsection*{As-built dimensions}

\begin{center}
\begin{tabular}{lcc}
\toprule
 & Bao's antenna (Y1) & Nick's antenna (Y2) \\
\midrule
$L$ initial cut (mm)              & 30.0  & 30.0  \\
$L$ after trim (mm)               & 26.0  & 26.7  \\
$W$ (mm)                          & 67.0  & 67.0  \\
Transmission line width (mm)      & 2.632 & 2.632 \\
Transmission line length (mm)     & \fb{TBD} & \fb{TBD} \\
Inset distance $R$ (mm)           & \fb{TBD} & \fb{TBD} \\
\bottomrule
\end{tabular}
\end{center}

\section{Initial Measurements and Impedance Matching}

\subsection*{Initial resonance and trimming}

With both antennas connected, $S_{11}$ (Bao) and $S_{22}$ (Nick) were swept from 1.5--2.7~GHz to find the initial resonance. The frequency with the lowest return-loss magnitude was treated as the operating frequency.

\begin{center}
\begin{tabular}{lcc}
\toprule
 & Bao (Y1, $S_{11}$) & Nick (Y2, $S_{22}$) \\
\midrule
Initial $L$ (mm)                    & 30.0 & 30.0 \\
Initial $f_\text{ant}$ (GHz)        & \fb{TBD} & \fb{TBD} \\
Trim direction                      & cut down & cut down \\
Amount trimmed (mm)                 & 4.0      & 3.3      \\
Post-trim $L$ (mm)                  & 26.0     & 26.7     \\
Post-trim $f_\text{ant}$ (GHz)      & $\approx 2.443$ & $\approx 2.443$ \\
\bottomrule
\end{tabular}
\end{center}

Both antennas started at $L = 30$~mm but resonated well below 2.443~GHz on the first sweep. About 4~mm was shaved off Bao's antenna and 3.3~mm off Nick's to bring each resonance up into the WiFi band. The 13\% trim on Bao's patch is larger than the simple $L = c/(2 f_c\sqrt{\varepsilon_r})$ formula predicts, which lines up with the well-known fringing-field effect: the radiating aperture extends about $0.5h$ past each end of the metal patch, so the electrical length is roughly $L + 2(0.5h) \approx L + 1.4$~mm longer than the physical length. The formula treats only the geometric $L$, so the as-cut patches always resonate a bit low and we trim from there.

\subsection*{Input impedance measurement}

At the post-trim operating frequency, the input impedance was read from the Smith chart marker.

\begin{center}
\begin{tabular}{lcc}
\toprule
 & Bao's antenna & Nick's antenna \\
\midrule
$Z_{IN\_0}$ from Smith chart ($\Omega$) & \fb{TBD} & \fb{TBD} \\
Return loss at $f_c$ (dB)               & \fb{TBD} & \fb{TBD} \\
\bottomrule
\end{tabular}
\end{center}

\subsection*{Inset feed calculation \texorpdfstring{$R$}{R}}

Moving the feed inward by distance $R$ scales the input impedance as
\[
Z_{IN\_R} = Z_{IN\_0}\cos^2\!\left(\frac{\pi R}{L}\right) \quad\Longrightarrow\quad R = \frac{L}{\pi}\cos^{-1}\!\sqrt{\frac{Z_{IN\_R}}{Z_{IN\_0}}}
\]
where $Z_{IN\_R} = 50~\Omega$ is the target and $Z_{IN\_0}$ is the measured edge-fed impedance.

\textbf{Calculated $R$ (plug in measured $L$ and $Z_{IN\_0}$):}
\begin{itemize}[leftmargin=*]
  \item Bao's antenna: $L = \fb{TBD},\ Z_{IN\_0} = \fb{TBD}\ \Omega \ \Rightarrow\ R = \fb{TBD~mm}$
  \item Nick's antenna: $L = \fb{TBD},\ Z_{IN\_0} = \fb{TBD}\ \Omega \ \Rightarrow\ R = \fb{TBD~mm}$
\end{itemize}

\subsection*{After inset cuts}

Per the lab procedure, about half of the calculated $R$ was cut from each side of the patch, then the VNA was monitored in Log-Mag mode and additional symmetric cuts made until return loss bottomed out.

\begin{center}
\begin{tabular}{lcc}
\toprule
 & Bao's antenna ($S_{11}$) & Nick's antenna ($S_{22}$) \\
\midrule
Final $R$ (mm)                     & \fb{TBD} & \fb{TBD} \\
Return loss at $f_c$ (dB)          & \fb{TBD} & \fb{TBD} \\
Return loss at 2.400~GHz (dB)      & \fb{TBD} & \fb{TBD} \\
Return loss at 2.487~GHz (dB)      & \fb{TBD} & \fb{TBD} \\
\bottomrule
\end{tabular}
\end{center}

\textbf{Target:} return loss below $-20$~dB at $f_c$, and at least $-9.5$~dB across the whole 2.400--2.487~GHz band.

\subsection*{VNA screen captures}

\fb{Insert VNA screen capture (deliverable \#1): both antennas connected, four-quadrant view of $S_{11}$/$S_{22}$ Log-Mag + Smith, with markers at 2.400, 2.443, and 2.487~GHz. Both $S_{11}$ and $S_{22}$ must read below $-20$~dB at the 2.443~GHz design frequency to satisfy the deliverable.}

\section{More Measurements and Polarization Observations}

\subsection*{Azimuth pattern}

The Smith chart quadrants were reconfigured to show $S_{12}$ and $S_{21}$ Log-Mag instead. The two antennas were placed facing each other, separated by a fixed distance set by the VNA cable lengths. One antenna was held stationary; the other was rotated in $15^\circ$ increments out to $90^\circ$. The aligned ($0^\circ$) measurement was taken as the 0~dB baseline, and the values below are the drop from that baseline.

\begin{center}
\begin{tabular}{cccc}
\toprule
Angle ($^\circ$) & $S_{21}$ (dB) & $S_{12}$ (dB) & Drop from 0~dB baseline (dB) \\
\midrule
0   & \fb{TBD} & \fb{TBD} & 0.0 \\
15  & \fb{TBD} & \fb{TBD} & \fb{TBD} \\
30  & \fb{TBD} & \fb{TBD} & \fb{TBD} \\
45  & \fb{TBD} & \fb{TBD} & \fb{TBD} \\
60  & \fb{TBD} & \fb{TBD} & \fb{TBD} \\
75  & \fb{TBD} & \fb{TBD} & \fb{TBD} \\
90  & \fb{TBD} & \fb{TBD} & \fb{TBD} \\
\bottomrule
\end{tabular}
\end{center}

Antenna separation distance: \fb{TBD~cm}.

\fb{Insert polar (azimuth) plot from Excel showing the radiation pattern across $\pm 90^\circ$. Use the 6+ data points above; assume symmetry across $0^\circ$.}

\subsection*{Cross-polarization test}

With both antennas facing each other, one was tipped on its side so that its $E$-field axis was rotated $90^\circ$ relative to the other. $S_{21}/S_{12}$ was recorded, then compared to the value with the two antennas facing entirely away from each other.

\begin{center}
\begin{tabular}{lcc}
\toprule
Configuration & $S_{21}$ (dB) & $S_{12}$ (dB) \\
\midrule
Both facing, one tipped $90^\circ$ (cross-pol)    & \fb{TBD} & \fb{TBD} \\
Both facing away from each other (back-to-back)   & \fb{TBD} & \fb{TBD} \\
\bottomrule
\end{tabular}
\end{center}

\fb{Discuss: is cross-pol lower than back-to-back, or higher? If lower, the patches are well-polarized and back-to-back coupling is dominated by reflections off the room. If higher, explain why (likely near-field coupling that doesn't respect polarization, or reflections off nearby surfaces).}

For the chosen separation, the far-field boundary is $r_{ff} = 2D^2/\lambda_0 \approx 7.31$~cm (using $D = W = 67$~mm and $\lambda_0 = 12.28$~cm at 2.443~GHz, derived above). Our antenna spacing of \fb{TBD~cm} is well past this threshold, so the link is operating in the far-field and the cross-pol result reflects true polarization mismatch rather than near-field coupling.

\section{Discussion and Conclusion}

\fb{Summarize: did the resonance land where the design calculation predicted (within a few percent)? How many trim cuts were needed to center it on 2.443~GHz? How close did the inset feed get to $-20$~dB at $f_c$ and $-9.5$~dB across the band? Comment on the azimuth plot shape (broad main lobe, $-3$~dB beamwidth around $\fb{TBD}^\circ$). Mention the cross-pol result and what it implies about how cleanly the patches are linearly polarized. Close with a brief note on which step was the most error-prone (probably the symmetric inset cuts) and how that affected the final match.}

\end{document}
